\section{Platform, Profiles, and Security}
\label{sec:platform}

\subsection{Account System}
\label{sec:accounts}

The Account System is the single place where the user manages all online
accounts and external service connections. Set up once, available everywhere,
with full visibility into which application accesses what. The design principle
is straightforward: no application ever directly handles OAuth flows, stores
tokens, or manages credentials. All credential logic lives in a dedicated
\emph{Account Daemon}.

The daemon serves two related but distinct purposes. First, it manages
user-facing online accounts such as Google, Microsoft, and Nextcloud, including
OAuth flows, token refresh, and the GNOME Online Accounts bridge. Second, it
manages the credentials for external service Connections configured in Settings
(Section~\ref{sec:connections}): GitHub, GitLab, n8n, OpenClaw, and any other
Connection the user creates. Both credential types live in the Linux Kernel
Keyring and are accessed exclusively through the daemon. No other component
reads credentials from the Keyring directly.

\emph{OAuth~2.0}~\cite{rfc6749} is the standard authorization framework used
by cloud services to grant third-party applications access to user data without
sharing the user's password. The flow works by redirecting the user to the
service's login page, where they authenticate and approve a specific set of
\emph{scopes} (declared permissions such as ``read email'' or ``write
calendar''), after which the service issues a short-lived \emph{access token}
and a long-lived \emph{refresh token}. The access token is sent with each API
request; the refresh token is used to obtain new access tokens when the current
one expires. An application that holds both tokens has persistent access to the
user's account; protecting them is therefore critical.

\paragraph{Architecture.}
Applications declare what they need; the daemon handles the rest. When an
application needs access to a service, it sends a request to the Account
Daemon, which checks the application's permission profile, retrieves or
derives the appropriate token, logs the access in the Audit Log, and returns
what the application needs. The application then communicates with the
provider's API directly using that token. The daemon is the sole keeper of
all credentials; tokens never live in application memory longer than
necessary.

\paragraph{Service-level permissions.}
Applications request access at the service level, not the OAuth scope level.
The user sees human-readable service names and never raw OAuth scope strings.
When an application requests account access for the first time, the user sees
exactly which service is being requested, grants or denies it, and the
decision is stored in the application's unified permission profile as described
in Section~\ref{sec:permissions}. Account permissions are part of the same
system as filesystem, network, and graph permissions.

\begin{lstlisting}[language=bash, caption={Service-level permission view (Settings)}]
Google Account (personal@example.com)
  Mail      -> allowed for: Thunderbird
  Calendar  -> allowed for: Calendar App, AI daemon (read-only)
  Drive     -> allowed for: File Manager
  Contacts  -> not granted to anyone
\end{lstlisting}

\paragraph{Multiple accounts and resolution order.}
When multiple accounts of the same type exist, the system resolves which to
use via a three-tier model. The system default applies globally; an Isolated
Desktop context can declare its own defaults, so that switching to the Work
desktop automatically switches the active account context; and individual
applications can be manually pointed at a specific account regardless of the
desktop context. Resolution proceeds from most specific to least specific:
per-app override, then desktop default, then system default.

\begin{lstlisting}[language=bash, caption={Account resolution configuration}]
# resolved in order: per-app override -> desktop default -> system default
[defaults]
google = "personal@example.com"

[desktop.work.defaults]
google = "work@example.com"

[app."org.mozilla.thunderbird"]
# always uses personal account regardless of active desktop
google = "personal@example.com"
\end{lstlisting}

\paragraph{Token security.}
Credentials never leave the Account Daemon as-is. Two paths are used
depending on provider support. For providers that support OAuth scope
restriction, the daemon issues a short-lived derived token scoped only to
what the application declared, with a default lifetime of five minutes. The
full OAuth token and refresh token never leave the daemon at all. For
providers that do not support downscoping, the full token is handed to the
application but with two compensating controls: every token handoff is logged
in the Audit Log, and tokens are actively revoked when the application process
exits. All credentials at rest are stored in the Linux Kernel Keyring~\cite{archwiki:keyring}: a kernel-managed storage
area for cryptographic keys, tokens, and passwords that lives in kernel memory
rather than userspace files, making it inaccessible to other processes and not
visible in \texttt{/proc} or on disk. The daemon
loads them into memory only when actively needed and discards them immediately
after.

\paragraph{Offline behavior.}
Two failure modes are handled distinctly. When the provider rejects the token
because the user revoked access or the account password changed, the daemon
emits a single system notification directing the user to Settings and does not
repeat it; applications receive a clear \texttt{AccountAuthRequired} error
code and surface a gentle in-app hint. When the network itself is unavailable,
the daemon serves cached data from the last successful sync, clearly marked
with the timestamp of the last update. A calendar application shows last-known
events labeled with their age rather than an empty state. When connectivity
is restored, the daemon refreshes silently in the background and notifies
applications of updated data via the Event Bus.

\paragraph{Connection credential management.}
External service Connections configured in Settings are stored as named
entries in the Linux Kernel Keyring under the \texttt{lunaris/connections/}
namespace. The Account Daemon is the sole process that reads and writes these
entries. When a component requests the credential for a Connection it provides
the Connection identifier and its own component identity; the daemon checks
whether the requesting component is listed in the Connection's \texttt{allowed}
field before returning anything. Unauthorized requests receive
\texttt{EACCES} with no credential data.

Token refresh for Connections follows a two-path model. For providers that
issue refresh tokens alongside access tokens, the daemon attempts a silent
refresh in the background when the access token approaches expiry. If the
refresh succeeds the new token is stored in the Keyring and the Connection
continues working without interruption. If the refresh fails, or if the
provider does not issue refresh tokens at all, the daemon marks the Connection
as \texttt{expired}, returns \texttt{token\_expired} to requesting components
rather than a stale or invalid token, and sends a single High-priority
notification to the user with a deep link to
\texttt{settings://connections/<id>}. The notification is sent once and not
repeated until the user interacts with it or the Connection is renewed.

The daemon also maintains the outbound webhook endpoint whitelist for the
Event Bus. Whenever a Connection with outbound events is created, modified,
or deleted, the daemon pushes the updated whitelist to the Event Bus over a
dedicated internal channel. The Event Bus never queries the daemon at send
time; it uses the cached whitelist to avoid any latency in the hot event path.

\paragraph{GNOME Online Accounts bridge.}
\emph{GNOME Online Accounts}~\cite{gnome:goa} (GOA) is a GNOME framework
that centralizes cloud service credentials (Google, Microsoft, Nextcloud, and
others) and exposes them to applications via a D-Bus interface.
Applications built for GNOME that use GOA expect a D-Bus
interface at \texttt{org.gnome.OnlineAccounts}. The system provides a
GOA-compatible bridge that implements this interface exactly but sources all
data from the Account Daemon internally. Evolution, GNOME Calendar, and other
GOA-dependent applications work without modification. Internally every request
goes through the Account Daemon, through the permission system, and through
the Audit Log. The bridge translates; it does not add logic.

\subsection{Profile System}
\label{sec:profiles}

A profile is a complete, isolated user environment backed by a separate Linux
system user. Each profile has its own home directory, its own Knowledge Graph,
its own account configuration, and its own set of running system daemons.
Isolation is enforced at the kernel level, not by convention.

The profile system replaces the concept of isolated desktops within a session.
Rather than trying to isolate individual desktops within a single session, the
entire session is the isolation boundary. Switching profiles is a session
switch, not a window toggle.

\paragraph{Personal profiles.}
The primary profile, created at system setup, cannot be deleted. It has no
restrictions: full access to all installed applications, all accounts, and all
system features. This is the default context the user logs into.

\paragraph{Work and custom profiles.}
A user-created or organization-deployed profile for a specific context such as
work, a project, or a secondary identity. Backed by a separate Linux user with
its own encrypted home image via \texttt{systemd-homed}. Virtual desktops,
accounts, Knowledge Graph, and application configurations are all independent
from the Personal profile. Multiple custom profiles can coexist. Each is a
full Linux user; switching to one is a standard session switch with saved and
restored window state.

\paragraph{Locked profiles.}
A profile deployed by an institution for controlled sessions: exams, regulated
work environments. Backed by a temporary Linux user created at session start
and deleted when the session ends. The user cannot switch to another profile
while a Locked session is active; this is enforced at the compositor level, not
by UI alone. After the session ends, the Linux user and all associated data are
removed. No data from the Locked session survives into the Personal profile.

\paragraph{Guest mode.}
A lightweight temporary context for short-term use by another person. Unlike
other profile types, Guest does not create a second Linux user or a second
compositor session. It runs as a sandboxed environment on top of the current
session: a separate namespace, an empty home directory, no access to the host
user's data. The host session is suspended in the background and resumed
instantly when Guest ends. All Guest data is discarded on exit.

\paragraph{Technical foundation.}
Each non-Guest profile is a Linux system user. Profile creation is
\texttt{useradd} plus \texttt{systemd-homed} home image setup; profile
switching is a standard Linux user session switch; profile deletion is
\texttt{userdel} plus home image removal. Isolation is kernel-enforced.
\texttt{systemd-homed} provides per-profile LUKS-encrypted home images. When
a profile is not active, its home image is unmounted and its encryption key
is discarded from memory. Each profile runs its own instance of all system
daemons: Graph Daemon, Account Daemon, Event Bus. There is no shared daemon
state between profiles.

\paragraph{Data transfer between profiles.}
By default profiles are fully sealed from each other. Controlled data transfer
is configured per profile pair in Settings, with independent directional
control over clipboard, drag-and-drop, and a shared folder:

\begin{lstlisting}[language=bash, caption={Cross-profile transfer configuration}]
[profile.work.transfer]
clipboard              = "in-only"    # "off" | "in-only" | "out-only" | "both"
drag_drop              = "in-only"
shared_folder          = "~/Transfer/Work"
shared_folder_permission = "read-write"  # or "read-only" | "off"
\end{lstlisting}

A dedicated \emph{Transfer Daemon} is the only process with simultaneous
access to both profile namespaces. All cross-profile data flow goes through
it; nothing bypasses it. Whenever data flows out of a profile, a
non-dismissable toast appears for three seconds identifying the source profile,
destination, and file name, with an Undo option. The transfer proceeds but can
be cancelled within the window. All transfers are recorded in the Audit Log
with source, destination, transfer type, and data description.

Locked profiles have transfer hardcoded to off in all directions. This cannot
be overridden by the user or by an organization's configuration package; it is
a system-level invariant with no configuration surface.

\paragraph{Signed profile packages.}
Profiles can be deployed via a \emph{Signed Profile Package}, a
\texttt{.lenv} file created and signed by an organization. The user installs
it themselves and retains full control over their device. The organization
gets a tamper-evident, versioned configuration. This is the BYOD model: the
organization controls their context, the user controls their device.

\begin{lstlisting}[language=bash, caption={Signed profile package format}]
[environment]
id        = "com.acme.work-profile"
name      = "Acme Work"
version   = "2.1.0"
publisher = "Acme GmbH"
type      = "persistent"    # or "locked" for exam/controlled sessions

[transfer]
clipboard     = "in-only"
drag_drop     = "in-only"
shared_folder = "off"

[network]
require_vpn    = "acme-vpn"
allow_outbound = ["*.acme.com", "*.office365.com"]
block_outbound = ["*"]

[apps]
required = ["com.acme.timetracker", "org.mozilla.firefox"]
blocked  = ["com.valve.steam"]

[graph]
ai_access = "off"
export    = "off"

# Project templates bundled with this profile.
# Each entry is a path to a .project file relative to the package root.
# These are written to the profile's home directory on install and
# registered as active projects in the Knowledge Graph.
[projects]
templates = [
    "projects/acme-onboarding.project",
    "projects/acme-infrastructure.project",
]
\end{lstlisting}

Project templates bundled in the \texttt{[projects]} block are installed into
the profile's home directory at \texttt{\textasciitilde/.projects/<template-name>/}
on first login and registered as active \texttt{Project} nodes in the profile's
Knowledge Graph. Each template is a standard \texttt{.project} TOML file
(Section~\ref{sec:projects}) and follows the same format; the organization
can pre-configure name, description, tags, linked accounts, AI scope, and
Focus Mode rules. Project templates do not contain user data and are not
synchronized back to the organization; they are a starting configuration
only. The user can modify or delete them after installation without affecting
the profile package itself.

The package is signed by the organization's private key. Key distribution uses
TOFU by default: on first install the user sees the publisher name and key
fingerprint and confirms trust. For managed onboarding, IT can send an
enrollment link so that the organization's key is imported before any package
arrives, eliminating the TOFU prompt entirely.

Before installation, a full summary of every policy and restriction is shown.
Nothing is hidden. There are certain things an organization cannot do through
a signed package, enforced technically rather than by policy: the organization
cannot read the user's Personal profile or its Knowledge Graph, cannot access
any data outside the profile's Linux user boundary, cannot monitor activity
in other profiles, cannot remotely wipe or modify the device, and cannot
prevent the user from uninstalling the package. If \texttt{notify\_on\_uninstall}
was declared and accepted at install time, the organization receives a
notification when the package is removed, but the removal proceeds regardless.

\subsection{Gaming and Windows Compatibility}
\label{sec:gaming}

Gaming and Windows application compatibility are first-class features, not
afterthoughts. The required components are pre-installed and pre-configured;
no manual setup is required.

\paragraph{Wine and Proton.}
\emph{Wine} provides compatibility for Windows applications on Linux.
\emph{Proton}~\cite{proton} is Valve's fork of Wine, optimized for gaming
and maintained with extensive patches from Valve's Linux gaming team; for gaming specifically
it is substantially better than vanilla Wine. Both are pre-installed alongside
\emph{DXVK}~\cite{dxvk}, which translates Direct3D 9, 10, and 11 API
calls to Vulkan. Direct3D is the Windows graphics API; most games targeting
Windows use it rather than Vulkan directly. By translating at the API boundary,
DXVK allows these games to run on Linux via the Vulkan driver.
\emph{vkd3d-proton}~\cite{vkd3d:proton} is Valve's fork of the Wine project's
vkd3d library; it translates Direct3D 12 (the modern Windows GPU API) to
Vulkan, with extensive optimizations for gaming workloads that upstream vkd3d
does not include. Winetricks is included for installing Windows runtime components
into Wine prefixes.

Steam is available in the Store and works out of the box with Proton enabled.
Steam manages its own Proton version per game; the system does not interfere
with Steam's internal Proton management.

Non-Steam Windows games and applications use Wine directly, managed through
the first-party \emph{Wine Manager} application. The underlying components are
shared with Steam's installation rather than duplicated. Wine Manager is a thin
UI over the existing tooling providing prefix creation and deletion, per-prefix
Winetricks access, launch shortcuts, and per-prefix Proton and Wine version
selection.

DXVK and vkd3d-proton require Vulkan. Mesa provides Vulkan support for AMD
and Intel GPUs by default. Nvidia requires the proprietary driver for Vulkan
support; the Store provides explicit guidance for Nvidia users during setup.

\paragraph{Wine theming.}
By default Wine applications render with a Windows Classic or Windows 7-era
appearance that is visually inconsistent with a modern desktop. A first-party
Wine theme is maintained as part of this project, generated from the same
design token system used for GTK and Qt theming. It updates automatically
when the system theme changes. Wine applications do not look native, but they
look consistent with the rest of the system. The Wine theme is applied
automatically to all Wine prefixes managed by the system; users who prefer
the default Windows appearance can opt out per prefix.

\paragraph{Performance tooling.}
\emph{Gamemode}~\cite{gamemode} by Feral Interactive is pre-installed and
enabled by default. It applies system-level performance optimizations when a
game is running: the CPU governor switches to performance mode, the I/O
scheduler is tuned, and kernel scheduler hints are applied. Proton handles the
Gamemode activation call automatically; no manual configuration is required.

\emph{MangoHud}~\cite{mangohud} is pre-installed as an in-game overlay showing
frames per second, GPU and CPU usage, temperatures, and frame timing. It is
disabled by default and toggled via a keyboard shortcut.

\paragraph{Knowledge Graph integration.}
Gaming activity is tracked in the Knowledge Graph like any other activity:
which games were played, when, and for how long; which Wine and Proton versions
were used per game; and performance data from MangoHud when enabled, linked to
sessions. This makes queries such as "how many hours did I play last week" or
"which Proton version worked best for this game" available locally without
relying on any external service.

\subsection{Security and Privacy}
\label{sec:security}

Security is not a feature added on top of this system. It is a foundational
design constraint that shapes every architectural decision.

\paragraph{Core principles.}
Five non-negotiables drive everything in this section. First, the Knowledge
Graph, Audit Log, and AI context are local by default and never leave the
machine without explicit user action; there is no opt-out dark pattern, the
default is that nothing goes anywhere. Second, the system collects no usage
data, crash reports, or any other information without explicit user consent.
Third, the system works fully offline and without registration; cloud features
are opt-in. Fourth, every component gets exactly the permissions it needs and
nothing more, applying to filesystem access, network access, graph access, and
syscalls alike. Fifth, security must not destroy usability: defaults are
chosen to protect without requiring constant interaction, because a system so
locked down that users disable its protections is less secure than a
well-calibrated one.

A sixth principle follows from the others: data minimization is itself a
security property. Data that does not exist cannot be stolen. The retention
policy defined in Section~\ref{sec:knowledge-graph} is not only a privacy
decision but an active security measure. Every decision to extend retention
windows is a deliberate expansion of attack surface and must be understood
as such.

\subsubsection{Threat Model}
\label{sec:threat-model}

Security measures without a named adversary are hard to evaluate. This section
states explicitly what this system is and is not designed to resist, so that
the measures described in the sections that follow can be judged against a
concrete baseline.

\paragraph{In scope.}
The primary adversary is a \emph{local attacker with user-level privileges}: a
malicious or compromised application running as the logged-in user that attempts
to exceed its declared permissions, exfiltrate graph data, escalate privileges,
or persist across sessions. The secondary adversary is an \emph{opportunistic
remote attacker} who exploits a network-reachable vulnerability in a running
service or a browser-class application. The tertiary adversary is a
\emph{physical attacker} with brief unsupervised access to a powered-off or
locked device. Managed-environment scenarios add a fourth adversary: a
\emph{malicious or over-reaching administrator} who attempts to access user
data beyond the boundaries described in Section~\ref{sec:managed}.

\paragraph{Out of scope.}
This system does not attempt to resist a \emph{root-level attacker}: an
adversary who already holds superuser privileges on the running system. Root
can modify kernel modules, bypass all userspace controls, and read any
unencrypted memory. Full Disk Encryption and TPM-sealed keys protect data at
rest against root on a different boot, not against root in the current session.
\emph{State-level adversaries} with hardware implant capabilities, supply-chain
compromise of the base distribution, or CPU microcode vulnerabilities are
similarly out of scope; no single OS can address these without hardware-level
attestation infrastructure outside the scope of this project.

\paragraph{Knowledge Graph as a high-value target.}
The Knowledge Graph introduces a risk that does not exist on a conventional
Linux system: a structured, queryable, cross-application timeline of user
activity. On a conventional system, an attacker who gains access to a user's
files reads unindexed data spread across many directories. An attacker who
gains access to the running Graph Daemon gains a pre-indexed, graph-traversable
record of everything the user has done. This is structurally more valuable than
a raw filesystem breach, and the security model of this system must account for
it explicitly.

The measures described in Section~\ref{sec:graph-security} reduce the
probability of such a breach. Section~\ref{sec:graph-daemon-compromise}
addresses what happens when prevention fails.

\subsubsection{Full Disk Encryption}
\label{sec:fde}

Full Disk Encryption means that everything stored on disk is encrypted at
rest. On Linux this is implemented via \emph{LUKS} on top of the
\emph{dm-crypt} kernel module. The entire partition is encrypted and the OS
decrypts it transparently at boot after credentials are provided.

FDE is enabled by default. A stolen or lost device is a realistic threat for
any user; without FDE, physical access means full access to all data. The
performance concern that historically argued against FDE no longer applies:
all modern x86 and ARM processors have hardware AES acceleration~\cite{intel:aes}
that makes encryption effectively free at current workloads.

The primary unlock mechanism is a passphrase entered at boot. On systems with
a \emph{TPM} (Trusted Platform Module)~\cite{archwiki:tpm}, a dedicated
security chip on the motherboard (or firmware-emulated since TPM~2.0) that
provides hardware-backed cryptographic operations, secure key storage, and
tamper-evident measurements of the boot process. The FDE key can be sealed to the TPM and released automatically
if the system boots in a known-good state. This allows passwordless boot on
trusted hardware while still protecting the drive if it is physically removed.
If the boot configuration changes due to tampering, the TPM refuses to release
the key and falls back to the passphrase. TPM-based unlock is opt-in;
passphrase is always the fallback.

Hibernate requires additional care because the entire RAM contents are written
to the swap partition. That swap partition must also be LUKS-encrypted, and
key management must allow the system to resume without manual intervention.
This is a known-solved problem on Linux but requires deliberate configuration
during the installation process.

\subsubsection{Secure Boot, Unified Kernel Images, and Measured Boot}
\label{sec:secureboot}

\emph{Secure Boot}~\cite{archwiki:secureboot} is a UEFI~\cite{uefi} firmware feature that verifies the cryptographic
signature of every piece of code executed during the boot process. If any
component is not signed by a trusted key, the firmware refuses to execute it.
The practical effect is that an attacker who modifies the bootloader or kernel
cannot boot the modified system without also possessing the signing key.

Both candidate base distributions support Secure Boot via a Microsoft-signed
shim bootloader. Custom kernel modules built for this project must be signed
with a Machine Owner Key enrolled in the firmware.

A traditional Linux boot chain consists of multiple separate files: the
bootloader, the kernel image, and the initramfs. A \emph{Unified Kernel Image}~\cite{archwiki:uki}
packages all three into a single signed EFI binary, reducing the attack
surface to a single component requiring verification. \texttt{systemd-boot}
supports UKI natively and the broader Linux ecosystem is moving in this
direction. This project targets UKI as the boot format.

\emph{Measured Boot} goes further than Secure Boot: every boot component is
cryptographically hashed and the hash is stored in the TPM's Platform
Configuration Registers. This creates a tamper-evident record of exactly what
ran during boot. If any component changes, even legitimately due to an update,
the PCR values change. Locally, the TPM can be configured to release the FDE
key only if the PCR values match an expected set, so that a modified bootloader
causes the TPM to refuse key release. In managed environments, a remote server
can request a signed attestation report from the TPM, and a device that fails
attestation can be denied network access.

\emph{Secure Boot} combined with \emph{Measured Boot} makes the Evil Maid
Attack significantly harder. An Evil Maid Attack is the scenario where an
attacker has brief physical access to a powered-off device and modifies the
bootloader to capture the user's FDE passphrase on the next boot. Without
the signing key, Secure Boot blocks the modified bootloader from executing
entirely.

\subsubsection{App Sandboxing}
\label{sec:sandboxing}

Every application runs in a restricted environment that limits what it can do
even if fully compromised. Three mechanisms operate simultaneously and
independently; bypassing one does not bypass the others.

\begin{table}[htbp]
\caption{Sandboxing mechanism layers}
\label{tab:sandboxing}
\begin{tabularx}{\linewidth}{lX}
\toprule
\textbf{Mechanism} & \textbf{Function} \\
\midrule
Linux Namespaces~\cite{archwiki:namespaces} & Changes what the application can see \\
AppArmor~\cite{archwiki:apparmor}           & Blocks actions on resources the application can see \\
seccomp~\cite{archwiki:seccomp}             & Blocks specific syscalls regardless of resource visibility \\
\bottomrule
\end{tabularx}
\end{table}

Linux Namespaces give each application its own restricted view of the system
before any permission check occurs. A Mount Namespace means the application
sees only its own data directory, explicitly granted paths, and a minimal set
of system paths; other applications' data directories do not exist from its
perspective. A Network Namespace means applications without declared network
access see no network interface at all; applications with declared network
access see only a filtered virtual interface. A PID Namespace means the
application can only see its own process tree.

AppArmor profiles are defined externally in the package and reviewed during
Store submission; an application cannot define or modify its own profile. A
text editor does not need \texttt{ptrace}, \texttt{mount}, or
\texttt{kexec}; a seccomp profile that excludes these syscalls causes the
kernel to kill the process immediately if a compromised application attempts
them.

IPC is restricted to the Event Bus, the Graph Daemon, and explicitly declared
inter-application communication channels. An application cannot open a socket
to another application without that channel being declared and approved.

\paragraph{Clipboard hygiene.}
The clipboard is an underappreciated data-leak vector. On a conventional
Wayland desktop, any application with a \texttt{wl\_data\_device} interface
can read the current clipboard contents without restriction; a password manager
and a text editor have equal access to whatever the user last copied. This
system changes that.

Every clipboard entry carries a sensitivity label. Labels are assigned in two
ways. Manual labeling is available from the clipboard history UI: the user can
mark an entry as sensitive, causing it to be hidden from clipboard history and
restricted from low-trust applications. Automatic labeling applies heuristics
at copy time: content that matches the pattern of an API key, an IBAN, a
private key, a password-manager field, or a credit card number is labeled
sensitive automatically without user interaction.

Applications declare a clipboard permission level in their manifest, analogous
to filesystem or network permissions. An application without the
\texttt{clipboard.read.sensitive} permission receives an empty string when it
requests clipboard contents that carry a sensitive label. The application is
not informed that content exists but was withheld; from its perspective the
clipboard is empty. This prevents a compromised application from silently
harvesting credentials that were copied from a password manager.

Clipboard entries labeled sensitive are excluded from clipboard history storage
entirely and are not written to the Knowledge Graph. They exist only in the
compositor's in-memory clipboard buffer and are cleared when the entry is
superseded or when the screen locks.

\subsubsection{Knowledge Graph Security}
\label{sec:graph-security}

The Knowledge Graph is the most sensitive component in the system. It
accumulates a detailed record of everything the user does and its access
control must be correspondingly rigorous.

No application has direct access to the Ladybug database files. All graph access
goes through the \emph{Graph Daemon}, a privileged system process that is the
sole reader and writer of the graph. Every other component communicates only
with the daemon over a defined IPC interface, and the daemon checks permissions
before executing any query.

\paragraph{Three axes of permission.}
Graph permissions are defined along three independent axes. The first axis is
node types: which categories of data the caller can read, such as
\texttt{File}, \texttt{App}, \texttt{NetworkConnection}, or \texttt{Session}.
A calendar application may read \texttt{Session} nodes but not
\texttt{NetworkConnection} nodes. The second axis is relations: which
connections between nodes the caller can traverse. An application permitted
to read \texttt{File} nodes may still be blocked from traversing the
\texttt{OPENED\_BY} relation that connects files to the applications that
opened them. The third axis is instances: whether the caller can see only the
nodes it generated itself or nodes generated by other applications as well.

\begin{lstlisting}[language=bash, caption={Graph permission entry}]
[app."com.example.notes"]
allow_read       = ["File.path", "File.last_modified", "Session.started_at"]
deny_read        = ["File.content_snippet", "NetworkConnection.*"]
allow_relations  = ["PART_OF_SESSION"]
instance_scope   = "self"   # only nodes this app generated
\end{lstlisting}

\paragraph{Two permission layers.}
Permissions are assigned in two layers. At install time, the package declares
the maximum graph permissions it can ever receive; the user reviews this
declaration as part of installation. At runtime, the application can request
permissions within its declared maximum, with a user prompt explaining what is
requested and why. A runtime request that exceeds the install-time declaration
is rejected outright. Some fields within permitted node types are additionally
restricted regardless of node-level permissions: file content snippets, for
example, require an explicit extra permission separate from general
\texttt{File} node access.

\paragraph{Ladybug file protection.}
The Ladybug database files are protected by three independent layers. At the
filesystem layer, the Graph Daemon runs as a dedicated system user
\texttt{graph-daemon} and the database files are owned exclusively by this
user with \texttt{600} permissions. A successful sandbox escape by an
application process running as the logged-in user does not grant access to
these files. At the AppArmor layer, a separate deny rule in every
application's AppArmor profile prohibits access to the graph directory even
if filesystem permissions were bypassed. At the encryption layer, the graph
directory lives inside a LUKS-encrypted loopback image that the daemon mounts
at startup using a key retrieved from the Kernel Keyring. When the daemon
stops, cleanly or via crash, the image is unmounted and the key is discarded
from memory. During any period when the Graph Daemon is not running, the Ladybug
files are LUKS-encrypted ciphertext unreadable even to root.

\paragraph{Cypher injection prevention.}
The Graph Daemon never accepts query strings with embedded values. All queries
must use named parameters, and the daemon runs a validator on every incoming
query that detects embedded values before execution.

\begin{lstlisting}[language=Rust, caption={Parameterized graph queries}]
// Rejected - string interpolation
let query = format!("MATCH (f:File {{name: '{}'}}) RETURN f", user_input);

// Required - parameterized
daemon.query(
    "MATCH (f:File {name: $name}) RETURN f",
    params!{ "name" => user_input }
)
\end{lstlisting}

Applications without \texttt{arbitrary\_query} permission in their manifest
cannot send free-form query strings at all; they select from pre-approved
query templates and supply only the parameters. Only explicitly trusted
components such as the AI daemon and system daemons receive
\texttt{arbitrary\_query} permission, and even they go through the
parameterization validator.

\paragraph{Capability tokens.}
Instead of looking up permissions in an ACL table on every query, applications
receive a cryptographically signed capability token at startup. The token
encodes the application's allowed scopes directly, and the daemon verifies the
\emph{HMAC} (Hash-based Message Authentication Code) signature on each query
without a database lookup: the verifier recomputes the HMAC and checks it
matches without needing to contact any authority. This is faster at high query
rates and more robust against a corrupted ACL table. The AI daemon can derive
restricted sub-tokens for specific query scopes without involving the daemon.
Token revocation when permissions change mid-session is handled via a small
in-memory revocation list that is cleared on daemon restart.

\paragraph{Token issuance and the Policy Store.}
Token issuance is separate from token verification and requires careful design:
if the component that signs tokens can be influenced to grant excessive scopes,
the entire permission model is undermined.

Tokens are signed by the Account Daemon at application startup. The permitted
scopes for each application are not decided by the Account Daemon itself;
instead, the daemon reads them from a \emph{Policy Store}, a static TOML file
owned by root with \texttt{644} permissions (world-readable, root-writeable
only). The Policy Store maps application IDs to their maximum permitted graph
scopes:

\begin{lstlisting}[language=bash, caption={Policy Store entry (read-only by Account Daemon)}]
[app."com.example.notes"]
graph_read  = ["File.path", "File.last_modified", "Session.started_at"]
graph_write = []
instance_scope = "self"
arbitrary_query = false
\end{lstlisting}

The Account Daemon can only issue a token whose scopes are a subset of what the
Policy Store declares for that application ID. It cannot grant scopes that are
absent from the Policy Store entry. A compromised Account Daemon can sign
fraudulent tokens, but only within the ceilings the Policy Store enforces.
The Graph Daemon independently reads the Policy Store and rejects any
token whose scopes exceed the registered maximum for that application. Scope
expansion requires a root-level file write, which is an observable, auditable
action. New applications receive a restrictive default scope at install time;
privilege elevation is explicit and logged.

\paragraph{Graph Daemon hardening.}
The daemon runs with the minimum privileges needed for its function, enforced
by a \texttt{systemd} service unit with \texttt{PrivateNetwork=true},
\texttt{MemoryDenyWriteExecute=true}, \texttt{PrivateTmp=true}, and a tightly
scoped \texttt{SystemCallFilter}. \texttt{PrivateNetwork=true} prevents a
compromised daemon from making outbound connections to exfiltrate data.
A \texttt{systemd} watchdog monitors the daemon; an unresponsive or crashed
daemon triggers a clean shutdown sequence that unmounts the LUKS image and
discards the key before restarting.

Clients verify that they are speaking to the legitimate Graph Daemon before
sending queries. At startup, the daemon signs a challenge with a private key
that only it holds; clients verify the signature against the public key
distributed with the system package. A connection that fails verification is
dropped, preventing a trojanized replacement daemon from intercepting graph
traffic.

\paragraph{Read/write process split.}
\label{sec:graph-daemon-split}
A single daemon process that both writes graph data and answers queries is a
larger attack surface than necessary. The Graph Daemon is therefore split into
two cooperating processes: \texttt{graph-daemon-write} accepts events from the
Event Bus and writes to Ladybug, and \texttt{graph-daemon-read} answers queries
from applications and daemons.

\texttt{graph-daemon-read} opens the Ladybug database in read-only mode. A
compromised read process cannot modify graph data. \texttt{graph-daemon-write}
does not expose a query interface to external callers; a compromised write
process cannot be queried for data it has written. The two processes communicate
only via a restricted internal channel to coordinate schema changes and to
propagate index updates. Both remain subject to the full systemd sandboxing
described above.

This split does not eliminate the risk of a compromised read daemon leaking
data it has legitimate access to. It does prevent an attacker who gains
code execution in one path from trivially pivoting to the other.

\paragraph{Compromise scenario and damage control.}
\label{sec:graph-daemon-compromise}
The measures above reduce the probability of a Graph Daemon compromise. They
do not make it impossible. An attacker who successfully executes arbitrary code
inside \texttt{graph-daemon-read} can issue arbitrary Cypher queries against
the open database handle and receive responses over the Unix socket, regardless
of the capability token the calling client presented. This is the honest
consequence of any architecture where a process must be running and reachable to
serve data: the running process is the trust boundary.

Several controls bound what an attacker can do from this position.
\texttt{PrivateNetwork=true} means the data cannot be sent over the network
from the daemon process itself; exfiltration requires additional steps through
another compromised component. The read-only Ladybug handle prevents writing or
deleting data. The Anomaly Detector monitors the Audit Log for query volume
spikes and unusual node-type access patterns; a compromised daemon issuing a
full-graph dump will generate anomaly alerts unless the attacker also
compromises the Audit Daemon, which is a separate process. The LUKS-encrypted
graph image means the data is protected at rest the moment the daemon stops.

What this system cannot offer is runtime confidentiality against a compromised
read daemon with an open database handle. No local architecture that provides
queryable history can offer this guarantee: it is a structural property of the
threat model, not a gap specific to this design. The correct response is
layered mitigation and honest disclosure to the user that the Knowledge Graph,
while protected against external breach, relies on the integrity of the
Graph Daemon process itself.

\paragraph{Rate limiting and query complexity.}
Two attack vectors are addressed: denial of service via query flooding, and
timing side-channels. Rate limiting is enforced per client identity using a
token bucket: normal applications are limited to 100 queries per second with
a burst limit of 200; the AI daemon has a higher sustained limit appropriate
to its workload; system daemons are unlimited. Exceeding the rate limit
generates an event for the Anomaly Detector, and repeated violations trigger
an alert.

Query complexity limits prevent expensive traversals from being used as denial
of service. Traversal depth is capped at 5 hops by default and never more than
10. Every query must include a \texttt{LIMIT} clause; queries without one are
rejected before execution. Queries are aborted after 500 milliseconds; the
client receives \texttt{QueryTimeout}. Response times have a small random
noise value added to make statistical timing attacks on the graph impractical.

\subsubsection{AI Security}
\label{sec:ai-security}

The AI daemon has broader graph access than most applications by necessity,
making it a higher-risk component with its own security model.

\paragraph{AI access levels.}
The user configures a global AI access level as the primary control over how
much the AI can see. Table~\ref{tab:ai-access} lists the five levels. Action
permissions, meaning what the AI can do rather than what it can read, are
configured separately. Reading data and taking actions are fundamentally
different risk categories and are therefore independently controlled. Suggest
only mode means the AI proposes actions and the user manually executes them.
Supervised mode means the AI executes actions but shows a preview with a
cancellation window before proceeding. Autonomous mode means the AI can act
within specific applications without per-action confirmation, and each
application must be individually enabled for this; autonomous mode is never a
global setting.

\begin{table}[htbp]
\caption{AI read access levels}
\label{tab:ai-access}
\begin{tabularx}{\linewidth}{llX}
\toprule
\textbf{Level} & \textbf{Name} & \textbf{What the AI can access} \\
\midrule
0 & Minimal       & Nothing from the graph; only what the user explicitly pastes \\
1 & Session-scoped & Only data from the current active session \\
2 & Project-scoped & All data associated with a user-selected project. Applied automatically by Focus Mode when a project is active (Section~\ref{sec:focus-mode}). \\
3 & Time-scoped   & Historical data up to a configurable lookback window \\
4 & Full read     & Everything the user has read access to \\
\bottomrule
\end{tabularx}
\end{table}

\paragraph{Prompt injection.}
Prompt injection is an attack where malicious instructions are embedded in
content the AI processes with the intent of manipulating its behavior: a PDF
might contain hidden text instructing the model to send files to an external
server. Current language models cannot reliably distinguish between content
they should process and instructions they should follow. This is a structural
limitation, not a fixable bug. The defense is therefore architectural.

Every piece of content the AI receives is tagged with its origin before
entering the prompt, so the model sees explicit markers distinguishing user
input, external content such as PDFs, and graph data. Before external content
reaches the main model, it passes through a small local model whose sole
function is to detect AI-directed instructions; this is a probabilistic first
pass that catches obvious injection attempts cheaply and locally.

As a hardcoded rule with no configuration surface, any action the AI wants to
take that was triggered by external content requires explicit user confirmation
regardless of the configured action permission level. An AI that read a
malicious PDF cannot use that as a basis to execute commands without the user
seeing exactly what it wants to do and why.

Documents are parsed in an isolated subprocess before their text content is
passed to the AI. The subprocess has no network access and no access to the
graph; only the extracted text leaves the sandbox, with no embedded scripts,
no active content, and no metadata that could carry instructions.

These layers substantially reduce the attack surface and raise the cost of a
successful injection, but they do not solve the underlying problem. Prompt
injection is an open, unsolved problem across the AI
industry~\cite{greshake2023prompt}. No current system can guarantee that a
sufficiently sophisticated injection embedded in external content will be
detected or neutralized. This system's position is to minimize exposure, make
AI-triggered actions visible and reversible, and be transparent about the
remaining risk rather than to claim a solution that does not exist.

\paragraph{AI network proxy.}
The AI daemon has no direct network access. All cloud API calls go through a
dedicated system proxy daemon that is the only component allowed to make
outbound connections to AI provider endpoints. The AI daemon constructs a
request, hands it to the proxy, and receives only the response. This prevents
a compromised AI daemon from independently exfiltrating data over arbitrary
network connections, ensures all AI-related network traffic is logged and
attributable, and allows the proxy to enforce that only approved provider
endpoints are contacted.

\subsubsection{Audit Log}
\label{sec:audit-log}

Every access to the Knowledge Graph, every AI action, and every permission
grant or denial is recorded in the Audit Log. This serves two purposes:
giving the user visibility into what is happening on their system, and
providing the data needed for behavioral anomaly detection.

The Audit Log is an append-only ledger. Entries cannot be modified or deleted;
only new entries can be added. Each entry is cryptographically linked to the
previous one via a hash chain, so tampering with any historical entry is
detectable by any reader that verifies the chain.

Two log tiers coexist. The Structural Log is always active and records
interaction metadata without content: timestamp, application identity, which
node types were queried, which relations traversed, result count, duration,
and status. The query string itself, result contents, and concrete node IDs
are never recorded in this tier. This is sufficient to detect suspicious
behavior without being sufficient to reconstruct what data was accessed.

The Forensic Log is opt-in and time-limited. When activated in Settings, it
additionally records the full query string, query parameters (not result
contents), and the calling process stack trace. Forensic Mode is always
limited to a maximum of 24 hours and then automatically deactivates. It is
shown prominently as active in Settings while enabled and cannot be activated
remotely by any application or administrator; only the user, locally.

A dedicated Audit Daemon is the sole writer to the log files. Other components
send audit events over a restricted IPC channel and cannot write directly to
the log. The Forensic Log is accessible exclusively to the user; no
administrator, no daemon, and no application can read it. The Structural Log
is readable by the user and by the Anomaly Detector for pattern analysis;
in managed environments, administrators receive aggregated anomaly alerts only,
never raw log contents.

\paragraph{Project-scoped export.}
The Audit Log supports filtering by project context for export. Settings
exposes a date-range export that can be scoped to a specific project: the
export includes all graph access records, AI actions, and permission events
that occurred during sessions where that project was active in Focus Mode.
The output is a structured JSON document suitable for time tracking,
client billing, or project handover. The Structural Log tier is used for
this export; the Forensic Log is never included in project exports. The
export is produced locally and handed to the user as a file; it is not
transmitted anywhere automatically.

\paragraph{Hash chain integrity and TPM anchoring.}
The hash chain makes retroactive tampering with individual log entries
detectable by any reader that verifies the chain. It does not, however, prevent
an attacker with early system access from deleting the log and constructing a
new chain from scratch over fabricated entries, because without an external
anchor the chain-start value is self-referential.

To close this gap, the HMAC key used to sign Audit Log entries is derived from
a secret sealed to the TPM against the measured boot state (TPM~2.0 PCR
sealing). The TPM releases the key only when the system booted in the exact
state recorded at seal time: the same bootloader, the same kernel, the same
initramfs. An attacker who wants to forge a clean audit chain must reproduce the
exact measured boot state that authorized the signing key, which requires the
Secure Boot signing key. The TPM sealing described here uses the same PCR
infrastructure already required for FDE key release in
Section~\ref{sec:fde}; no additional hardware is needed.

On systems without a TPM the HMAC key falls back to a software-derived secret
stored in the Kernel Keyring. This preserves the chain-integrity guarantee
against post-hoc entry modification but not against early-access forgery. The
Security section of Settings indicates which protection level is active; this
difference is disclosed to the user.

\subsubsection{Anomaly Detection}
\label{sec:anomaly}

The Anomaly Detector is a system daemon that reads the Audit Log and
identifies unusual patterns. It is not a traditional antivirus and does not
match file signatures or scan for known malware. It is a behavioral analysis
tool that establishes a baseline of normal behavior on this specific machine
and flags deviations.

Signature-based detection requires a continuously updated database of known
threats and cannot detect novel attacks. Because this system already maintains
a rich picture of normal behavior via the Knowledge Graph and Audit Log,
behavioral detection is the natural fit: the data needed to define normal
behavior already exists.

The Anomaly Detector watches for patterns including an application suddenly
querying graph node types it has never queried before; query volume from an
application spiking above its historical baseline; an AI action executed
without a preceding user interaction in the same session; a network connection
from a process that has never previously made network connections; a new eBPF
program being loaded outside system update windows; and permission escalation
requests at unusual times.

Alerts are surfaced via the notification system with enough context for the
user to make a decision: which application, what it did, why it is unusual,
and options to investigate or dismiss. The detector does not automatically
block anything; it informs. Automatic blocking based on anomaly scores is an
opt-in feature for managed environments where an administrator configures the
policy. The Anomaly Detector runs locally and alert data does not leave the
machine unless the user or administrator explicitly configures remote alert
forwarding.

\subsubsection{Network Security}
\label{sec:network-security}

On a conventional Linux system, any process can open a network connection to
anywhere by default. This is the wrong default. On this system, outbound
network access is denied unless explicitly permitted. Every application that
needs network access declares it in its package manifest by protocol,
destination, and direction; connection attempts outside the declaration are
blocked at the kernel level via eBPF-based network filtering and logged as
security events.

VPN enforcement can be configured in managed environments such that certain
network destinations are reachable only when a VPN is active. \emph{WireGuard}~\cite{wireguard} is the preferred VPN protocol. WireGuard is
a modern VPN implemented directly in the Linux kernel (merged in 5.6), using
Curve25519 for key exchange and ChaCha20-Poly1305 for encryption. Its design
is deliberately minimal: the entire implementation is under 5,000 lines of
kernel code, compared to tens of thousands for OpenVPN, making it easier to
audit and maintain. Key management is simpler than IPsec: each peer has a
static key pair, and peers are added by exchanging public keys.
\emph{OpenVPN}~\cite{openvpn} support is provided for compatibility with
existing enterprise infrastructure.

\paragraph{Tor and i2P.}
Two anonymity networks are supported as opt-in features configurable in the
Advanced section of Network Settings. Both are pre-installed and require no
manual setup beyond enabling them; the complexity that typically prevents
adoption is handled by the system.

\emph{Tor}~\cite{tor} (The Onion Router) is an anonymity network that routes
traffic through a series of volunteer-operated relays, encrypting it in layers
so that no single relay knows both the origin and the destination of a
connection. Each relay knows only its predecessor and successor in the chain;
the full path is visible to no single party. Tor is widely used by journalists,
activists, and privacy-conscious users to circumvent surveillance and
censorship, and is operated by the non-profit Tor Project.

When Tor routing is enabled in Advanced Settings, the system configures a
transparent proxy using \texttt{nftables}: all outbound TCP traffic is
redirected through the local \texttt{tor} daemon, and DNS resolution is also
routed through Tor to prevent DNS leaks. UDP traffic cannot traverse the Tor
network and is blocked entirely while Tor routing is active rather than allowed
to bypass it, which would leak information about the user's network activity.
This is the correct behavior and is disclosed clearly in the Settings UI.
Tor routing does not provide anonymity against an adversary who controls the
entry and exit relays simultaneously, against JavaScript fingerprinting in
the browser, or against application-level identifiers such as cookies and
logged-in accounts. The Settings UI states these limitations plainly; Tor
routing is presented as a tool, not a guarantee.

\emph{i2P}~\cite{i2p} (Invisible Internet Project) is a separate anonymity
network optimized for internal services rather than general internet access.
Unlike Tor, which is primarily an anonymizing proxy for the regular internet,
i2P is a self-contained network where participants host and consume services
within the network itself. Websites within i2P (called \emph{eepsites}) are
reachable only from within i2P; external internet access through i2P is
possible but secondary to its design intent. i2P uses a distributed hash table
for routing and encrypts traffic in layers analogously to Tor, but with a
different threat model and different traffic analysis resistance properties.

The system uses \emph{i2pd}~\cite{i2pd}, a C++ implementation of the i2P
router. The original i2P router is written in Java and carries a significant
runtime overhead; i2pd provides the same network participation and tunnel
routing in a fraction of the memory footprint, with no Java dependency. i2pd
runs as a system service managed by systemd and is controlled through a
first-party Tauri management UI that exposes tunnel configuration, bandwidth
limits, and network participation status without requiring the user to edit
configuration files manually.

Both Tor and i2P are entirely disabled by default and leave no listening
services or background processes active until explicitly enabled. Enabling
either requires navigating to Advanced Network Settings, not a quick-access
toggle, reflecting that these tools require the user to understand their
purpose and limitations.

\subsubsection{Physical Security}
\label{sec:physical}

Full Disk Encryption handles the most common physical threat, but physical
access threats go further.

A Cold Boot Attack exploits the fact that DRAM does not lose its contents
instantly when power is removed. At low temperatures, an attacker who obtains
a running or recently suspended device could potentially read the FDE key from
memory. Two mitigations apply: the kernel is configured to overwrite RAM
contents before suspend-to-disk, and AMD SME and Intel TME are hardware
features available on modern CPUs that encrypt the entire RAM contents with a
key that lives inside the CPU die and never leaves it. Even if an attacker
physically removes the RAM modules, they read only ciphertext. Both are
enabled by default where hardware supports them.

\emph{USBGuard}~\cite{usbguard} enforces a policy on USB device connections. When the screen
is locked, no new USB devices are permitted to become active; only devices
connected before the screen locked remain accessible. This blocks BadUSB
attacks, where a specially crafted USB device presents itself as a keyboard
and injects keystrokes. Enabled by default; specific devices such as hardware
security keys can be whitelisted in Settings.

\emph{IOMMU}~\cite{archwiki:iommu} (Input-Output Memory Management Unit)
restricts which memory regions each hardware device can access via Direct
Memory Access (DMA). With IOMMU enabled, a malicious Thunderbolt or
PCIe device can only access the memory regions the OS has explicitly mapped
for it, not arbitrary system RAM. Enabled by default via kernel parameters.
For Thunderbolt specifically, the default security level requires newly
connected devices to be explicitly authorized by the user before they are
activated; previously authorized devices are remembered.

TPM Remote Attestation allows a device to cryptographically prove to a remote
server that it booted in a known-good state by sending a signed report of its
measured boot PCR values. In a managed deployment this can gate network access:
a device that fails attestation is denied VPN access until an administrator
investigates. This is inspired by the \emph{BeyondCorp}~\cite{beyondcorp} model
(Google's zero-trust network architecture, where no device or user is trusted
by default based on network location alone) and represents a meaningful
differentiator for enterprise deployments.

\subsubsection{Memory Safety}
\label{sec:memory-safety}

Memory safety bugs are the most common source of exploitable vulnerabilities
in system software and are structural problems with C and C++ that persist even
with experienced developers. All system components developed for this project
are written in Rust. The Rust compiler enforces memory safety at compile time
through its ownership and borrow checker; an entire class of bugs that would
be possible in C simply cannot compile in Rust. This is a compile-time
guarantee, not a runtime check with overhead.

Third-party applications cannot be rewritten. Sandboxing limits what a
compromised application can do; hardware-level mitigations provide a further
layer. \emph{Intel CET}~\cite{intel:cet} (Control-flow Enforcement Technology)
adds a hardware-enforced shadow stack: a second, protected copy of return
addresses that the CPU maintains separately from the regular call stack. On
every function return, the CPU compares the return address from the regular
stack against the shadow stack; a mismatch causes an immediate fault. This
defeats the most common class of control-flow hijacking, return-oriented
programming (ROP), where an attacker overwrites saved return addresses to chain
together existing code fragments. CET is supported by the Linux kernel since
version 6.6 and is enabled by default for new processes on supported hardware
with no application changes required. For any C or C++ components compiled as part of this project,
Clang's \emph{CFI}~\cite{clang:cfi} (Control Flow Integrity) instrumentation
is applied. CFI adds compile-time instrumentation that validates at runtime
that every indirect function call and virtual method dispatch targets only a
function with the expected type signature; an attacker who corrupts a function
pointer to point to an arbitrary location will cause an immediate abort rather
than code execution.

\subsubsection{Multi-User Isolation}
\label{sec:multiuser}

When multiple users share a system, their isolation goes deeper than
traditional Unix file permissions because the Knowledge Graph, AI
configuration, and Audit Log introduce sensitive state beyond the home
directory. Each user has a separate Ladybug instance, a separate Graph Daemon
process running under their own system user, and a separate Audit Log. These
are not shared resources with access control layered on top; they are
separate instances that other users' processes cannot reach.

Home directories are managed via \texttt{systemd-homed}: each user's home
directory is stored as a LUKS-encrypted portable image. When a user logs out,
their home image is unmounted and the encryption key is discarded from memory.
No other user, including root, can access the home image without the user's
credentials. The Knowledge Graph lives inside the home directory and is
therefore protected by this encryption as a structural consequence rather than
an additional policy.

In a multi-user or managed environment, an administrator can install and remove
system-wide software, configure network policies, set password and screen lock
policies, view aggregated anomaly alerts, manage USBGuard whitelists, and
manage Secure Boot keys. The administrator explicitly cannot read any user's
Knowledge Graph, access any user's home directory contents, view any user's
Audit Log, or intercept any user's AI sessions. These limitations are enforced
by \texttt{systemd-homed} and the per-user Graph Daemon architecture, not by
policy alone. The administrator sees that something unusual happened via
anomaly alerts; they do not see what.

\subsection{Managed Environments}
\label{sec:managed}

For enterprise deployments where this system runs on many machines administered
centrally, additional infrastructure is required beyond what individual Signed
Profile Packages provide. This section describes the managed environment
architecture for organizations that need fleet-level control.

\paragraph{Policy Engine.}
A system daemon receives signed policy packages from a central management
server and applies them locally. The daemon verifies the cryptographic
signature of every package before applying any policy; unsigned or incorrectly
signed packages are rejected outright. Policies can control permitted and
blocked applications, network access rules beyond per-app declarations, AI
provider restrictions including complete disabling of cloud AI, the USBGuard
device whitelist, screen lock timeout and password complexity requirements,
update scheduling so that updates are applied only during configured maintenance
windows, and a minimum audit log retention period for compliance purposes.

\paragraph{LDAP and Active Directory.}
Most enterprises already have an identity management system, typically
Microsoft Active Directory or an LDAP-compatible directory. Integration
with these systems is a baseline requirement for enterprise adoption.
\emph{LDAP} (Lightweight Directory Access Protocol~\cite{rfc4511}) is the
standard protocol for querying and modifying directory services such as
Microsoft Active Directory and OpenLDAP. It stores organizational information:
users, groups, computers, and their attributes. \emph{Kerberos}~\cite{rfc4120}
is the authentication protocol layered on top: it issues time-limited tickets
that prove identity without sending passwords over the network; Active Directory
uses Kerberos as its primary authentication mechanism.

User authentication, group membership, and basic policy delivery are handled
via these standard protocols, which existing enterprise infrastructure already
supports without proprietary agents on the management server side.

\paragraph{MDM-style management.}
Beyond LDAP, a modern MDM-style management layer provides richer capabilities:
device enrollment, remote policy push, health status reporting, and remote
wipe. This is a longer-term goal; LDAP compatibility is the foundation and
MDM-style management is built on top of it once the base is stable.

\paragraph{Remote attestation as network gate.}
TPM Remote Attestation, described in Section~\ref{sec:physical}, can be used
as a network access gate in managed deployments. The workflow is:

\begin{enumerate}
\item The device boots and generates a TPM attestation report containing signed
  PCR values.
\item The device connects to the management server's attestation endpoint.
\item The server verifies the report against the expected baseline for that
  device's configuration.
\item If verification passes, the device receives its VPN credentials and
  network access.
\item If verification fails, the device is quarantined pending administrator
  review.
\end{enumerate}

This ensures that a device with a modified bootloader, disabled Secure Boot, or
an outdated security policy cannot silently access company resources. It is
inspired by the BeyondCorp model~\cite{beyondcorp} and is a meaningful
differentiator for enterprise deployments.

\paragraph{GDPR and compliance tension.}
In managed environments there is an inherent conflict between user privacy
rights and organizational compliance requirements. GDPR grants users rights
over their personal data; many compliance frameworks in financial services,
healthcare, and legal contexts require organizations to retain logs and audit
trails. This system does not resolve the conflict, because it is a legal and
organizational question rather than a technical one.

What the system does is make both sides technically possible. Users retain
full visibility into and control over their own Knowledge Graph and personal
Audit Log. Organizations can configure a minimum retention period for
compliance-relevant audit events, which do not include Knowledge Graph
contents. The separation between personal graph data, which is private to
the user, and system audit events, which may be subject to organizational
retention requirements, is maintained architecturally rather than by policy.
Organizations deploying this system should seek legal counsel to define
appropriate policies. The system provides the technical primitives; the
legal framework is out of scope.
